\subsection{序群的增同态}

设 G 和 \(G'\) 是两个序群；在从G的底层加法群到 \(G'\) 的底层加法群的同态中，考虑增映射是有意义的，即满足 \(x \leqslant y\) 蕴含 \(f(x) \leqslant f(y)\) 的那些映射。由于关系 \(f(y - x) = f(y) - f(x)\) ，由此可知，从 G 到 \(G'\) 的增同态的特征在于：在这样的同态下，G 的正元素的像是 \(G'\) 的正元素；若 P（相应地， \(P'\) ）表示 G（相应地，G'）的正元素集合，这可以写成 \(f(P) \subset P'\) 。显然，子群 G 到序群 G' 的典范单射，以及序群的乘积到其一个因子上的投影，都是增同态。

从序群 G 到序群 \(G'\) 上的一个同构（VI，第2页）f 是从 G 到 \(G'\) 上的双射同态，使得 f 与其逆同态都是增的，换言之， \(f(P) = P'\) 。

可能发生这样的情况：从 G 的底层群到 \(G'\) 的底层群的同构是增的，而逆同构却不是增的。例如，当 \(G = G'\) 时，若 f 是 G 上的恒等映射，且 \(P \subset P'\) 但 \(P \neq P'\) ，就是这种情形。因此在 Z 中，我们可以取 \(P'\) 为（通常意义的）正整数集合，而 P 为正偶整数集合。

(DIV) 设 K 是有理函数域 \(\mathbf{F}_{2}(\mathbf{X})\) （在 2 阶域 \(F_{2}\) 上）。相对于环 \(F_{2}[X] = A'\) 与环 \(F_{2}[X^{2}, X^{3}] = A\) 的整除关系在 \(K^{*}\) 上定义了两个不同的序群结构，它们满足 \(A \subset A'\) （它们确实是序群结构，因为 1 是 A 中唯一的单位，也是 \(A'\) 中唯一的单位）。

